\subsection{H2: is deletion a cheaper way to find untested code?}

No. Under equal 20-minute wall-clock budgets, deletion mutants yielded \HTwoRatio{} times as many activated survivors per CPU-minute as mutmut's classical operators (95\% CI \HTwoLo{} to \HTwoHi{}). The ratio was above 1 in \HTwoAbove{} of \HTwoN{} projects (Table~\ref{tab:h2}), so H2 is \textbf{\HTwoSupported}. This is the clearest negative result of the study, and the opposite of what we expected from the Java literature, which presents extreme mutation as the economical option~\cite{vera2018descartes,deng2013deletion}.

Two measurements explain it. First, deletion mutants survive less often: \HTwoDelSurvShare{} of completed deletion mutants survived, against \HTwoClsSurvShare{} of classical ones. Removing a whole function body, or a call, is usually a large behavioural change that some test notices. Second, a deletion mutant is not cheaper to run. The median CPU cost per mutant is almost the same (\HTwoDelCpuMedian{}\,s against \HTwoClsCpuMedian{}\,s), but the mean is higher for deletion (\HTwoDelCpuPerMutant{}\,s against \HTwoClsCpuPerMutant{}\,s). Deletion mutants are covered by more tests (median scope \HTwoDelScopeMedian{} against \HTwoClsScopeMedian{}), break more of them when killed (median \HTwoDelKillersMedian{} failing tests against \HTwoClsKillersMedian{}), and time out about twice as often (\HTwoDelTimeoutShare{} against \HTwoClsTimeoutShare{}); we did not investigate the individual timeouts.

The deletion population is also small. It ran to completion inside the budget in \HTwoExhausted{} of \HTwoN{} projects, using \HTwoDelCpuMin{} CPU-minutes in total against \HTwoClsCpuMin{} for the classical runs, which ran out of budget in five projects. Completing the whole deletion population costs about a quarter of what the classical budget spent, and gives \HTwoDelSurv{} survivors against \HTwoClsSurv{}.

Survivors per CPU-minute is a measure of yield, not of value. A surviving classical mutant is often a small change of a constant or an operator, a population known to contain many equivalent mutants~\cite{papadakis2019advances}, while a surviving deletion says that a whole function or call can disappear unnoticed. The pre-registered usefulness review was designed to compare the two on that axis. At the time of writing, \HTwoReviewed{} of the 50 sampled survivors have been labelled, so that comparison is pending (real-gap share: classical \HTwoClsRealGap{}, deletion \HTwoDelRealGap{}). Blinding the reviewer to strategy turned out to be impractical, since a deletion patch is recognisable on sight.

\begin{table}[t]
\centering\small
\begin{tabular}{lrrrrrr}
\toprule
Project & Classical surv. & Classical CPU min & Deletion surv. & Deletion CPU min & Ratio & Deletion exhausted \\
\midrule
attrs & 300 & 158.49 & 9 & 15.89 & 0.30 & yes \\
cachetools & 193 & 9.86 & 14 & 9.94 & 0.07 & yes \\
click & 462 & 159.36 & 69 & 59.31 & 0.40 & yes \\
dateutil & 377 & 149.87 & 31 & 39.86 & 0.31 & yes \\
marshmallow & 318 & 62.91 & 14 & 9.59 & 0.29 & yes \\
more-itertools & 201 & 155.65 & 4 & 22.26 & 0.14 & yes \\
tomlkit & 566 & 158.63 & 85 & 44.40 & 0.54 & yes \\
toolz & 161 & 14.88 & 2 & 1.26 & 0.15 & yes \\
\bottomrule
\end{tabular}
\caption{H2: activated, valid survivors and CPU minutes under equal 20-minute wall-clock budgets. \emph{Ratio}: deletion survivors per CPU-minute divided by classical survivors per CPU-minute. \emph{Deletion exhausted}: every eligible deletion mutant ran within the budget.}
\label{tab:h2}
\end{table}

